\scanpage{62}
\uline{Corollaire} Soit $C$ un cube de dim 3. Alors
\[
\operatorname{Aut}^+(C)\xrightarrow{\sim}\operatorname{Aut}_{\mathrm{ens}}(\underbrace{\Phi_0(C)/\sigma}_{\text{ens. de card. 4}})
\]

\uline{Corollaire} Le foncteur $C\longmapsto\Phi_0(C)/\sigma$ est une équivalence de la catégorie des 3-cubes \inserted{\uline{orientés}} dans \deleted{\unsure{de dim 3|dans $\to$}} la cat. des ens. de card. 4 (en prenant comme morphismes les iso \ldots)

\uline{Cube (combinatoire) gauche}
\[
C=C(B,\sigma)\subset E=E_C=E(B,\sigma)
\]
cube. Considérons le quotient par la relation d'antipodisme, mais sur
\[
\partial\unsure{\hat C|\widetilde C|C^*|\overleftarrow C}=\coprod_{F\in\Phi^{**}(\deleted{\unsure{E}}C)}\overline F
\]
\[
\partial\hat C/\sigma\simeq E^*/\mathbb R^*=P(E)\qquad
\begin{aligned}&\text{espace projectif réel}\\&\text{de dim }n-1\end{aligned}
\]

La décomposition de $\partial\hat C$ en facettes donne une décomposition (où $C\overset{p}{\unsure{-|\to}}\partial C$ passage au quotient)
\[
\widetilde C\deleted{\unsure{-|=}}=\partial\hat C/\sigma
=\coprod_{F\in\Phi^{**}}p(F)=\coprod_{F\in\Phi^{**}}p(\overline F)
\]

NB Pour tte $F\in\Phi^{**}$, $\overline F\to P(E)$ est injectif

On a \deleted{$p(\overline F\cap\overline F')=p(\overline F)\cap p(\overline F')$ si $\overline F\cap\overline F'=$}
\[
p(\overline F)\cap p(\overline F')=
\begin{cases}
\varnothing\quad\text{si }\overline F\cap\overline F'=\varnothing\quad\text{et }\overline F\cap\sigma\overline F'=\varnothing\\
[p(\overline F)\cap p(\overline F')]\unsure{\cup|\mathbin{\dot\cup}}
[p(\overline F)\unsure{\cap\sigma|\cap|\overline{\cap}}p(\overline F')]\quad\text{sinon}\\
\deleted{\unsure{p(\overline F)\cap p(\sigma\overline F')\quad\text{si }\overline F\cap\sigma\overline F'\ne\varnothing}}
\end{cases}
\]
[\uline{NB} les relations $\overline F\cap\overline F'\ne\varnothing$ et $\overline F\cap\sigma\overline F'\ne\varnothing$ ne s'excluent pas]

\deleted{\unsure{i.e. si} $B\to I$ et $B',B''$ deux sections partielles}

et $p(\overline F)\cap p(\overline F')$ n'est pas néc. connexe : il est connexe sss \deleted{\unsure{une seule}} on n'a pas ces 2 relations à la fois)
% archive sheet number: 61
